\chapter{Una máquina de neuronas vivas}
\chaptermark{Una máquina de neuronas vivas}
\label{sec:livingmachine}

\section{Un banco de pruebas con el circuito a la vista}

En el capítulo~\ref{sec:debugging} depuramos una máquina de la que no tenemos el esquema: la sonda oye mil neuronas de ochenta mil millones y todo lo demás hay que adivinarlo. Un ingeniero en esa situación haría algo sencillo: montaría un pequeño trozo del circuito en una placa de pruebas, donde se ve cada componente. Con neuronas también se puede hacer.

Las neuronas de la corteza se separan y se cultivan sobre un chip con una matriz de electrodos, de unas pocas decenas a decenas de miles. Al cabo de unas semanas las células extienden sus prolongaciones y se conectan en una red de miles o cientos de miles de neuronas, sobre todo \nt{GLUT} y una fracción menor de \nt{GABA} que depende de la preparación: en los cultivos de hipocampo es de alrededor del $6\%$ de las neuronas~\cite{Benson1994}. Cada electrodo sabe escuchar, como la sonda del capítulo~\ref{sec:debugging}, y también inyectar corriente, como una señal de prueba. La segunda variante del banco son los \emph{organoides}: bolitas tridimensionales de tejido nervioso de alrededor de un milímetro, cultivadas a partir de células madre humanas~\cite{Lancaster2013}. En estos bancos las redes vivas funcionan durante meses; hay plataformas donde se pueden hacer experimentos con organoides a distancia, por la red, durante más de cien días seguidos~\cite{Jordan2024}. Esta área se llama \emph{computación biológica} o \emph{inteligencia organoide}~\cite{Smirnova2023}.

El banco difiere del cerebro en dos cosas importantes. Es diminuto: tiene cien mil veces menos neuronas. Y no tiene canales de difusión: en un cultivo de corteza no hay neuronas \nt{DOPA}, \nt{SERO}, \nt{NORA} ni \nt{ACET}. Es una máquina con bus de datos y control de flujo, pero sin registros de control. Veamos de qué es capaz.

\section{Lo que la red hace por sí sola}

Dejado en paz, el cultivo no calla ni produce un ruido uniforme. De vez en cuando casi toda la red se enciende a la vez, en una \emph{ráfaga de red}, y vuelve a apagarse; los intervalos entre ráfagas van de un segundo a varios minutos, según el cultivo~\cite{Wagenaar2006}. Para un ingeniero es un cuadro conocido: un amplificador con fuerte realimentación positiva y sin carga se convierte en un oscilador.

El mecanismo ya lo conocemos. Las fuertes conexiones mutuas entre neuronas \nt{GLUT} desencadenan una avalancha, como en el capítulo~\ref{sec:gaba} cuando no se cumplen las condiciones de la proposición~\ref{prop:ei}. La avalancha agota rápido la reserva de neurotransmisor en las sinapsis~\eqref{eq:depression}, y la red enmudece. La reserva se recupera con constante de tiempo $\tau_{\text{rec}}$, y cuando se ha recuperado la fracción $x^*$ suficiente para una nueva avalancha, llega una nueva ráfaga. El intervalo entre ráfagas es
\begin{empheq}[box=\keybox]{equation}
T_{\text{ráfaga}} \;\approx\; \tau_{\text{rec}}\,\ln\frac{1}{1-x^*} .
\label{eq:burstinterval}
\end{empheq}
Con $\tau_{\text{rec}}=0.8\,\text{s}$ del capítulo~\ref{sec:code} y $x^*=0.9$ salen unos dos segundos; con $x^*=0.99$, unos cuatro. Es una estimación del modelo con el $\tau_{\text{rec}}$ del libro, y cae en el extremo corto de los intervalos medidos. La medición directa en microcultivos da una recuperación de la reserva de neurotransmisor tras la ráfaga en decenas de segundos~\cite{CohenSegal2011}, es decir, allí $\tau_{\text{rec}}$ es mayor; con ese $\tau_{\text{rec}}$ la fórmula~\eqref{eq:burstinterval} da decenas de segundos y minutos, como en los cultivos lentos. Es un oscilador de relajación, pariente del generador de ritmo del capítulo~\ref{sec:muscles}: allí lo recargaba la fatiga de los grupos; aquí, la reserva de neurotransmisor.

Las ráfagas son lo que la red hace cuando no tiene nada que hacer. En el cerebro vivo se ve un cuadro parecido en el sueño profundo (capítulo~\ref{sec:sleep}) y en el cerebro en desarrollo, antes de que le lleguen entradas de verdad. El parecido sugiere una idea, pero que los mecanismos sean los mismos es una hipótesis.

\section{Cómo enseñar a la red sin un maestro dopaminérgico}

Como en el cultivo no hay \nt{DOPA}, la señal de “mejor o peor” hay que darla nosotros, con los electrodos. Los experimentos muestran que la red del banco sí cambia sus respuestas, y lo más interesante en ellos es \emph{con qué} se la enseña.

\emph{Un maestro que se calla.} Se estimula la red por un electrodo una vez cada uno a tres segundos hasta que en otro electrodo aparece la respuesta deseada $50\pm10\ms$ después del estímulo. En cuanto aparece, la estimulación se apaga durante cinco minutos. Tras varios ciclos así, la respuesta deseada aparece enseguida ante el estímulo~\cite{Shahaf2001}. Aquí la recompensa es el silencio: el estímulo desestabiliza la red, y al cesar deja a la red en el estado que dio la respuesta correcta. Es el descenso por el gradiente mediante desviaciones aleatorias del capítulo~\ref{sec:dofa} (proposición~\ref{prop:gradient}), donde el papel del error lo desempeña la propia sacudida: se sacude hasta que sale bien.

\emph{Una red que juega al tenis.} En un experimento en un banco con una matriz densa de electrodos, alrededor de un millón de neuronas se conectaron a un videojuego de tenis con una sola raqueta~\cite{Kagan2022}. La posición de la pelota se entregaba con corriente a los electrodos “sensoriales”: dónde está la pelota, con un código de lugar; a qué distancia, con la frecuencia. La actividad en dos regiones “motoras” movía la raqueta hacia arriba o hacia abajo. La regla de aprendizaje era inusual. Si la raqueta devolvía la pelota, la red recibía un estímulo breve y \emph{predecible}; si fallaba, varios segundos de corriente aleatoria \emph{impredecible}. En veinte minutos de juego las series de pelotas devueltas se alargaban. Un aumento significativo dentro de la sesión solo se dio con la retroalimentación completa; con silencio en lugar del estímulo, el cultivo al final de la sesión también jugaba mejor que sin retroalimentación, pero con un efecto menor, y no hubo aprendizaje estable de una sesión a otra a lo largo de varios días. El análisis detallado de este experimento está en el capítulo~\ref{sec:dishbrain}.

Los autores explicaron el resultado con la hipótesis de que la red actúa de modo que su entrada sea predecible~\cite{Friston2010}. Es la misma idea que el ahorro de espigas del capítulo~\ref{sec:code} y la predicción del fotograma del capítulo~\ref{sec:recognition}: la máquina gasta menos cuando la entrada es la esperada. Pero tanto el experimento como, sobre todo, las palabras de los autores sobre la “sintiencia” del cultivo provocaron duras críticas: el efecto es pequeño y se puede explicar de forma más sencilla~\cite{Balci2023}. La valoración honesta es esta: el cultivo en el banco cambia su comportamiento bajo la retroalimentación, y eso está medido; que aprenda precisamente a predecir su entrada es una hipótesis.

\emph{Una red que conduce un cuerpo.} De modo parecido, se conectó un cultivo a un sencillo “cuerpo” virtual y se le enseñó a moverse hacia un objetivo eligiendo el patrón espaciotemporal de los estímulos de entrenamiento~\cite{Bakkum2008}. En ninguno de estos experimentos hay dopamina: el maestro es un patrón de corriente que elige el ingeniero. Qué cambia cuando el maestro pasa a ser un programa o la propia dopamina se ve en las secciones~\ref{sec:dishdopamine}--\ref{sec:cartpole}.

\begin{figure}[t]
\centering
\begin{tikzpicture}[>=Stealth,
  blk/.style={draw,very thick,rounded corners=2pt,align=center,font=\footnotesize,minimum height=0.8cm,fill=blue!5}]
% (a) closed loop
\node[font=\small] at (1.5,4.3) {(a) lazo cerrado};
\node[blk,text width=2.6cm] (G) at (1.5,3.3) {juego: pelota y raqueta};
\draw[very thick,fill=orange!10,rounded corners=3pt] (0.5,0.2) rectangle (2.5,1.6);
\foreach \x in {0.7,1.0,...,2.35}{\foreach \y in {0.4,0.7,1.0,1.3}{\fill[gray!60] (\x,\y) circle (0.04);}}
\draw[->,thick,blue!60!black] (G.west) -| (-0.5,0.9) -- (0.5,0.9);
\node[font=\scriptsize,blue!60!black,anchor=east,align=right] at (-0.6,2.1) {pelota\\$\to$ corriente};
\draw[->,thick,red!70!black] (2.5,0.9) -- (3.5,0.9) |- (G.east);
\node[font=\scriptsize,red!70!black,anchor=west] at (3.6,2.1) {espigas $\to$ raqueta};
\node[font=\scriptsize] at (1.5,-0.2) {neuronas sobre la matriz de electrodos};
\node[font=\scriptsize,align=center] at (1.5,-0.85) {acierto: estímulo predecible\\fallo: corriente aleatoria};
% (b) reservoir
\begin{scope}[xshift=7.9cm]
\node[font=\small] at (1.6,4.3) {(b) reservorio};
\node[blk,text width=1.1cm] (X) at (-0.4,1.0) {entrada\\$u(t)$};
\draw[very thick,fill=orange!10] (1.6,1.0) circle (0.8);
\foreach \a/\r in {20/0.35,80/0.5,150/0.3,210/0.55,280/0.4,330/0.2,110/0.15}{\fill[red!60!black] ({1.6+\r*cos(\a)},{1.0+\r*sin(\a)}) circle (0.05);}
\node[font=\scriptsize] at (1.6,-0.2) {organoide: sin maestro};
\node[blk,text width=1.8cm,fill=green!8] (W) at (4.0,1.0) {lectura\\$W_{\text{sal}}$};
\draw[->,thick] (X) -- (0.8,1.0);
\draw[->,thick] (2.4,1.0) -- node[above,font=\scriptsize] {$\mathbf r(t)$} (W);
\draw[->,thick] (W) -- (4.0,2.3) node[above,font=\scriptsize] {$y(t)$};
\node[font=\scriptsize,green!40!black] at (4.0,0.3) {aprende con maestro};
\end{scope}
\end{tikzpicture}
\caption{Dos maneras de hacer que una red viva calcule. (a) Lazo cerrado: la posición de la pelota se entrega como corriente, las espigas de las regiones motoras mueven la raqueta, y el estímulo predecible o aleatorio tras el golpe hace de maestro. (b) Reservorio~\eqref{eq:reservoir}: al organoide no se le da señal de enseñanza; descompone la entrada en muchas respuestas no lineales con memoria, y por el error aprende una sola lectura lineal externa. Las conexiones del propio organoide también cambian, sin maestro.}
\label{fig:livingmachine}
\end{figure}

\section{Un reservorio vivo}

Hay otra manera de usar una red viva: no enseñarle nada. En ingeniería se llama \emph{computación de reservorio}~\cite{Maass2002,JaegerHaas2004}. Se toma un sistema no lineal con memoria ya hecho, cualquiera, siempre que su respuesta a la entrada sea rica y dependa del pasado reciente. La entrada $u(t)$ lo excita, se obtiene un vector de respuestas $\mathbf r(t)$, y solo se entrena la lectura lineal
\begin{empheq}[box=\keybox]{equation}
y(t) \;=\; W_{\text{sal}}\,\mathbf r(t) ,
\label{eq:reservoir}
\end{empheq}
cuyos pesos se ajustan con la regla de mínimos cuadrados~\eqref{eq:lms} del capítulo~\ref{sec:motorlearning}. El reservorio hace lo que hacían las pequeñas neuronas \nt{GLUT} del capítulo~\ref{sec:motorlearning}: descompone la entrada en un vector largo en el que las clases buscadas se vuelven separables por un plano (capítulo~\ref{sec:dendrites}).

Un organoide sobre una matriz de electrodos sirvió como reservorio de este tipo. Sonidos del habla, entregados como patrones de corriente, se distinguieron por hablante tras entrenar una sola lectura lineal, con una exactitud de alrededor del $78\%$~\cite{Cai2023}. Esa exactitud del $78\%$ es la cifra que dan los propios autores y que repite la prensa. El resultado es útil, pero importa entender dónde está la “inteligencia”: el organoide aporta no linealidad y memoria que se desvanece, y el aprendizaje por el error se hace fuera, en silicio (fig.~\ref{fig:livingmachine}). Con todo, el organoide no queda inalterado: según los autores, durante el entrenamiento cambió su propia conectividad funcional, y lo llaman aprendizaje no supervisado en el propio organoide~\cite{Cai2023}. Qué parte de la exactitud se debe a ese cambio y qué parte a la lectura no se ha separado.

\section{Qué nos dice el banco sobre la máquina}

\emph{Energía.} Según la estimación~\eqref{eq:energy}, una espiga cuesta unos $10^{-10}$~J. Un cultivo de un millón de neuronas a una espiga por segundo gasta en transmitir señales
\begin{empheq}[box=\keybox]{equation}
P \;\approx\; 10^{6}\times1\,\text{s}^{-1}\times10^{-10}\,\text{J} \;=\; 10^{-4}\,\text{W},
\label{eq:dishpower}
\end{empheq}
una décima de milivatio. La electrónica que estimula, registra y procesa esas espigas gasta órdenes de magnitud más. Como en el capítulo~\ref{sec:debugging}, el cuello de botella no es el cálculo, sino la comunicación con él.

\emph{Las piezas que faltan.} El banco muestra cuánto sabe hacer el bus \nt{GLUT} con el control de flujo \nt{GABA} sin registros de control: autoorganizarse, generar ráfagas, cambiar sus respuestas bajo la retroalimentación y servir como buen reservorio. Lo que sin ellos no sabe hacer es decidir por sí mismo qué cuenta como éxito. En el banco esa señal la da el ingeniero; en el cerebro, como vimos en los capítulos~\ref{sec:dofa}--\ref{sec:acet}, los canales de difusión.

\emph{Ética.} Los organoides se cultivan a partir de células humanas, y cuanto más complejos se vuelven, más serio es el problema de si ese tejido puede sentir algo. La mirada del ingeniero sugiere una respuesta parcial: el dolor del capítulo~\ref{sec:pain} no es la señal de un solo sensor, sino todo un sistema de alarma con compuertas, control de volumen y canales de difusión, que en el banco no existen. Pero esto es un razonamiento, no una prueba, y el propio campo propone hacer la evaluación ética junto con los experimentos, desde el principio~\cite{Smirnova2023}.

\begin{keyblock}
\begin{proposition}[Una red viva en el banco]
\label{prop:livingmachine}
Una red de neuronas \nt{GLUT} y \nt{GABA} sobre una matriz de electrodos genera por sí sola ráfagas~\eqref{eq:burstinterval}, cambia sus respuestas bajo la retroalimentación y puede servir como reservorio~\eqref{eq:reservoir} para una lectura entrenada desde fuera. Todo esto está medido. Que una red así aprenda según alguna regla general, por ejemplo a predecir su entrada, es una hipótesis, y las palabras grandilocuentes sobre su “sintiencia” no las acepta la mayoría de los especialistas.
\end{proposition}
\end{keyblock}

\section{Dopamina con un destello de luz}
\label{sec:dishdopamine}

El único experimento directo que conocemos en el que se dio dopamina a un cultivo como maestro se hizo en una plataforma de organoides a la que los investigadores se conectan a distancia~\cite{Jordan2024}. Al líquido que baña los organoides se le añadió dopamina en una “jaula” fotosensible: la molécula enjaulada no actúa sobre los receptores. La concentración de dopamina enjaulada es de $1$~mM. Cuando hay que recompensar a la red, se la ilumina con luz ultravioleta: la jaula se rompe y la dopamina sale al volumen.

El lazo funciona así. El mismo estímulo se aplica una y otra vez; si provocó más espigas que antes, se enciende el destello y se libera dopamina. A lo largo de $240$ repeticiones el número de espigas evocadas creció en conjunto. Los propios autores escriben que con un solo experimento así no se puede concluir que el aumento se deba a la dopamina~\cite{Jordan2024}.

Para un ingeniero es el primer intento de montar en una placa de cultivo la regla de tres factores~\eqref{eq:threefactor}: la actividad del emisor y del destinatario deja una traza, y una señal común decide si se consolida el cambio. Pero aquí la señal solo puede ser positiva: hay recompensa o no la hay, y el montaje no produce un error negativo $\delta<0$. Además, la dopamina se dispersa y se retira despacio, como la concentración en~\eqref{eq:lowpass}, así que los destellos frecuentes pueden simplemente elevar su nivel general. Para distinguir el aprendizaje de esto hacen falta dos controles: el mismo número de destellos en momentos aleatorios, sin relación con la respuesta de la red, y los mismos destellos sin dopamina enjaulada. Y hace falta una comprobación tras el descanso: si el cambio se mantiene cuando la dopamina ya se ha ido.

\section{La dopamina enseña al silicio}
\label{sec:biohybrid}

En el experimento inverso, las células vivas enseñan a la electrónica~\cite{Keene2020}. Células que liberan dopamina se cultivan en un microcanal sobre un elemento neuromórfico orgánico: un transistor cuya conductancia de canal hace de peso sináptico. La dopamina que sale de las células se oxida en el electrodo del elemento, y eso cambia la conductancia. El dispositivo reproduce también el mecanismo de retirada del neurotransmisor de la hendidura, de modo que el peso no solo se puede cambiar de forma duradera, sino también devolver a su valor.

En términos de circuito es un integrador con fugas con entrada química: el peso acumula el flujo de dopamina, y la “retirada” fija lo rápido que lo olvida; es el mismo circuito RC que~\eqref{eq:lowpass}. Lo principal aquí es la dirección de la conexión. La sonda del capítulo~\ref{sec:debugging} no ve los registros de difusión, porque son químicos. Este elemento lee precisamente un registro químico y lo convierte en un peso eléctrico. Para las interfaces cerebro-computadora es un camino hacia un sensor que oiga no solo el bus de datos, sino también los registros de control.

\section{Organoides con sus propias neuronas dopaminérgicas}
\label{sec:midbrainorg}

A partir de células madre humanas se saben cultivar organoides parecidos a la región del cerebro donde están las neuronas \nt{DOPA}. Contienen neuronas dopaminérgicas funcionales que liberan dopamina~\cite{Jo2016}. Estos organoides se cultivan sobre todo para estudiar enfermedades. No hemos encontrado experimentos en los que su dopamina sirviera de maestro para otra red.

Para una máquina de neuronas vivas, esta es la pieza que falta. El banco del capítulo~\ref{sec:livingmachine} es un bus de datos con control de flujo sin registros de control; aquí la fuente de la señal de difusión ya existe. Falta el crítico: una red que calcule el error de predicción~\eqref{eq:ctd} y gobierne estas neuronas. Conectar un organoide de corteza con uno de estos organoides de modo que la dopamina se libere según el error es un problema abierto, y por ahora es una hipótesis, no un experimento.

\section{Un organoide equilibra un péndulo sin dopamina}
\label{sec:cartpole}

El experimento reciente más completo prescinde por completo de la dopamina~\cite{Robbins2026}. Se conectaron organoides de corteza de ratón en lazo cerrado a la tarea del “péndulo sobre un carro”: la actividad del organoide mueve el carro, y la posición del péndulo se devuelve como estimulación. Entre intentos, el organoide recibe breves estímulos de enseñanza de alta frecuencia. Cuáles exactamente lo elige un programa de aprendizaje por refuerzo.

Los resultados están medidos:
\begin{itemize}
\item en la mayoría de los organoides, las señales de enseñanza elegidas por el programa dieron mejor resultado que señales elegidas al azar o que la ausencia de señal;
\item tras $45$ minutos de descanso la mejora no se conservaba;
\item el bloqueo de los receptores de glutamato de ambos tipos, AMPA y NMDA, anulaba la mejora.
\end{itemize}

El último punto dice dónde vive lo aprendido: en las sinapsis \nt{GLUT}, y a través del receptor que en la sección~\ref{sec:weightstore} funciona como detector de coincidencias entre entrada y salida. El segundo punto muestra lo que falta: hay cambio rápido, pero no consolidación, como el aprendiz rápido del capítulo~\ref{sec:motorlearning} sin el lento. Y el papel del maestro ha cambiado: al ingeniero que elige el patrón de corriente lo ha sustituido un programa, pero el maestro sigue estando fuera de la placa.

Entre los experimentos anteriores de aprendizaje de cultivos sobre matrices de electrodos tampoco hemos encontrado ninguno que usara dopamina: la señal de enseñanza fue siempre la estimulación eléctrica.

\section{Quién enseña al cultivo}

\begin{table}[t]
\centering
\footnotesize
\setlength{\tabcolsep}{4pt}
\renewcommand{\arraystretch}{1.15}
\begin{tabular}{>{\raggedright}p{2.7cm}>{\raggedright}p{2.5cm}>{\raggedright}p{3.2cm}>{\raggedright\arraybackslash}p{4.4cm}}
\toprule
Experimento & Quién enseña & Señal de enseñanza & Qué se sabe \\
\midrule
parar la estimulación ante la respuesta deseada & el ingeniero & fin de la estimulación & la respuesta se fija; medido \\
DishBrain (cap.~\ref{sec:dishbrain}) & el ingeniero & corriente predecible o aleatoria & aumento significativo de peloteos en la sesión solo con realimentación completa, con silencio el efecto es menor; medido; inestable de una sesión a otra; la causa se discute \\
péndulo sobre un carro & programa de aprendizaje por refuerzo & estímulos de alta frecuencia elegidos por el programa & mejor que los aleatorios, pero no se conserva tras el descanso; medido \\
dopamina por destello & la regla “más espigas = recompensa” & dopamina liberada por la luz & aumento de espigas: observación; el papel de la dopamina no está demostrado \\
sinapsis biohíbrida & células vivas & dopamina de las células & cambio duradero y retorno del peso del elemento; medido \\
organoides con neuronas \nt{DOPA} & --- & --- & la fuente de dopamina existe; no se ha usado como maestro \\
\bottomrule
\end{tabular}
\caption{Quién enseña a las neuronas vivas sobre un chip, y con qué.}
\label{tab:dishteachers}
\end{table}

En todos los experimentos en los que aprende el propio cultivo, el maestro está por ahora fuera (tabla~\ref{tab:dishteachers}): el ingeniero, un programa o una regla fijada por el ingeniero. La dopamina ha entrado en la placa una sola vez, y su papel aún está por demostrar. Montar en la placa un verdadero canal de difusión, con crítico, error y traza como en el capítulo~\ref{sec:dofa}, es el siguiente paso, que nadie ha dado todavía.


\section{Resumen}

\begin{table}[t]
\centering
\footnotesize
\setlength{\tabcolsep}{4pt}
\renewcommand{\arraystretch}{1.15}
\begin{tabular}{>{\raggedright}p{2.8cm}>{\raggedright}p{4.2cm}>{\raggedright\arraybackslash}p{5.8cm}}
\toprule
Experimento & Qué hace la red viva & Qué se sabe \\
\midrule
red sin entrada & ráfagas con intervalos de un segundo a minutos~\eqref{eq:burstinterval} & medido; el mecanismo por agotamiento sináptico es el modelo aceptado \\
maestro que se calla & la respuesta deseada se fija cuando se apaga la estimulación & medido \\
juego de tenis & las series de pelotas devueltas se alargan; un aumento significativo en la sesión, solo con retroalimentación completa & el cambio de comportamiento está medido; el aprendizaje de una sesión a otra es inestable; el tamaño del efecto y la explicación por predicción de la entrada son discutibles \\
cuerpo virtual & movimiento hacia el objetivo con un patrón de estímulos elegido & medido \\
reservorio & no linealidad y memoria para la lectura~\eqref{eq:reservoir} & medido; el aprendizaje por el error ocurre fuera, en silicio; las conexiones del organoide también cambian \\
energía & unos $10^{-4}$~W por millón de neuronas~\eqref{eq:dishpower} & estimación según~\eqref{eq:energy} \\
\bottomrule
\end{tabular}
\caption{Lo que mostraron las máquinas de neuronas vivas.}
\label{tab:livingmachine}
\end{table}

Una máquina de neuronas vivas (tabla~\ref{tab:livingmachine}) es un banco de pruebas en el que se ve qué saben hacer el bus de datos y el control de flujo sin registros de control. Saben hacer mucho, pero no saben decidir por sí mismos qué está bien. En el banco ese papel lo desempeña el ingeniero con su patrón de corriente; en el cerebro, los pocos cables de difusión a los que está dedicada la parte central del libro.

\section{Ejercicios}

\begin{exercise}[\lvlA]
\label{ex:21:1}
\emph{Intervalo entre ráfagas.} La ráfaga agota la reserva hasta cero; luego $\tau_{\text{rec}}\,\dd x/\dd t=1-x$, y la nueva ráfaga empieza cuando $x=x^*$; $\tau_{\text{rec}}=0.8$~s. (a)~Halle el intervalo para $x^*=0.9$ y $0.99$. (b)~El umbral $x^*$ varía al azar en $\pm0.005$ en torno a $0.99$. ¿En qué rango caen los intervalos?
\end{exercise}

\begin{exercise}[\lvlA]
\label{ex:21:2}
\emph{Energía del banco.} (a)~¿Qué potencia da la estimación~\eqref{eq:dishpower} para todo el cerebro ($8.6\cdot10^{10}$ neuronas)? (b)~El banco registra $1024$ electrodos a $20$~kHz con $10$~bit por muestra. ¿Cuántas veces más cara que la potencia del propio cultivo es la transmisión de este flujo a $1$~nJ por bit?
\end{exercise}

\begin{exercise}[\lvlB]
\label{ex:21:3}
\emph{Diseño: suprimir las ráfagas.} Una estimulación escasa por muchos electrodos hace que cada sinapsis libere neurotransmisor con frecuencia $r_b$; la reserva cumple $\tau_{\text{rec}}\,\dd x/\dd t=1-x-Ur_b\tau_{\text{rec}}x$, $U=0.5$, $\tau_{\text{rec}}=0.8$~s. ¿Con qué $r_b$ mínima la reserva no llega a $x^*=0.9$; $0.99$, y cuánto se debilita entonces la sinapsis?
\end{exercise}

\begin{exercise}[\lvlB]
\label{ex:21:4}
\emph{La memoria del reservorio no supera el número de neuronas.} Un reservorio con $N$ salidas $\mathbf r(t)$ recibe una entrada blanca $u(t)$ de media nula y varianza unidad; $m_k$ es el mayor cuadrado de la correlación de la lectura $\mathbf w_k\cdot\mathbf r(t)$ con $u(t-k)$. Demuestre que la capacidad de memoria $\mathrm{MC}=\sum_{k\ge0}m_k\le N$.
\end{exercise}

\begin{exercise}[\lvlB]
\label{ex:21:5}
\emph{Simulación: un reservorio en silicio.} Reservorio $\mathbf r(t+1)=\tanh(W\mathbf r(t)+\mathbf v\,u(t)+\mathbf c)$, $N=30$, $W$ aleatoria con radio espectral $0.9$, $v_i\sim U(-0.5,0.5)$, $c_i\sim U(-0.2,0.2)$, $u\sim U(-1,1)$; la lectura es una regresión de cresta. Halle la capacidad de memoria del ejercicio~\ref{ex:21:4} con $\tanh$ y sin él. ¿Qué reservorio recuerda más?
\end{exercise}

\begin{exercise}[\lvlB]
\label{ex:21:6}
\emph{Lectura para un reservorio vivo.} Un organoide da $1024$ canales y $P=240$ registros de entrenamiento de dos clases. (a)~¿Cuántos rasgos $n$ se pueden conservar para que, según la fórmula del ejercicio~\ref{ex:19:3}, un etiquetado aleatorio sea separable con probabilidad no mayor del $1\%$? (b)~¿Cuántas sigmas por encima del azar está una exactitud del $78\%$ en $100$ registros de prueba?
\end{exercise}

\begin{exercise}[\lvlC]
\label{ex:21:7}
\emph{Investigación: un maestro sin canal de difusión.} Proponga un protocolo de estimulación que implemente en el banco la regla de tres factores del capítulo~\ref{sec:dofa} a través de $8$--$1024$ electrodos, y un control con la lectura congelada que distinga el aprendizaje de pesos de un desplazamiento del estado de la red. ¿Qué resultado refutaría la hipótesis del aprendizaje?
\end{exercise}
